\textbf{COLUNI 2024 -- Questão 06}\par\medskip A obra \textit{Terra de mil povos}, de Kaká Werá, \textit{A carta pras Icamiabas}, de \textit{Macunaíma} e a música \textit{Índios}, da banda Legião em Urbana trazem como ponto comum a representação do indígena brasileiro. Considerando uma leitura comparativa das obras, é CORRETO afirmar que, o ponto de vista destas em relação ao indígena é:\begin{enumerate}
\item semelhante, no que se refere à representação do indígena como personagem principal das obras.
\item semelhante, no que se refere à narrativa do desprestígio histórico de como o indígena era tratado.
\item semelhante, no que refere à representação de três exemplos de literatura produzida por índios.
\item semelhante, no que se refere à construção dos enredos das obras na perspectiva do indígena.\end{enumerate}
